\documentclass[10pt,a4paper]{amsart}
\usepackage{amsmath,amssymb,amsthm,graphicx}
\usepackage{verbatim}
\newenvironment{reference}{\comment}{\endcomment}
%\newenvironment{reference}{}{}
\newenvironment{slogan}{\comment}{\endcomment}
\newenvironment{history}{\comment}{\endcomment}

% For commutative diagrams we use Xy-pic
\usepackage[all]{xy}

% We use 2cell for 2-commutative diagrams.
\xyoption{2cell}
% Theorem environments.
%
\theoremstyle{plain}
\newtheorem{theorem}[subsection]{Theorem}
\newtheorem{proposition}[subsection]{Proposition}
\newtheorem{lemma}[subsection]{Lemma}

\theoremstyle{definition}
\newtheorem{definition}[subsection]{Definition}
\newtheorem{example}[subsection]{Example}
\newtheorem{exercise}[subsection]{Exercise}
\newtheorem{situation}[subsection]{Situation}

\theoremstyle{remark}
\newtheorem{remark}[subsection]{Remark}
\newtheorem{remarks}[subsection]{Remarks}

\numberwithin{equation}{subsection}

% Macros
%
\def\lim{\mathop{\mathrm{lim}}\nolimits}
\def\colim{\mathop{\mathrm{colim}}\nolimits}
\def\Spec{\mathop{\mathrm{Spec}}}
\def\Hom{\mathop{\mathrm{Hom}}\nolimits}
\def\Ext{\mathop{\mathrm{Ext}}\nolimits}
\def\SheafHom{\mathop{\mathcal{H}\!\mathit{om}}\nolimits}
\def\SheafExt{\mathop{\mathcal{E}\!\mathit{xt}}\nolimits}
\def\Sch{\mathit{Sch}}
\def\Mor{\mathop{\mathrm{Mor}}\nolimits}
\def\Ob{\mathop{\mathrm{Ob}}\nolimits}
\def\Sh{\mathop{\mathit{Sh}}\nolimits}
\def\NL{\mathop{N\!L}\nolimits}
\def\CH{\mathop{\mathrm{CH}}\nolimits}
\def\proetale{{pro\text{-}\acute{e}tale}}
\def\etale{{\acute{e}tale}}
\def\QCoh{\mathit{QCoh}}
\def\Ker{\mathop{\mathrm{Ker}}}
\def\Im{\mathop{\mathrm{Im}}}
\def\Coker{\mathop{\mathrm{Coker}}}
\def\Coim{\mathop{\mathrm{Coim}}}

% Boxtimes
%
\DeclareMathSymbol{\boxtimes}{\mathbin}{AMSa}{"02}

%
% Macros for moduli stacks/spaces
%
\def\QCohstack{\mathcal{QC}\!\mathit{oh}}
\def\Cohstack{\mathcal{C}\!\mathit{oh}}
\def\Spacesstack{\mathcal{S}\!\mathit{paces}}
\def\Quotfunctor{\mathrm{Quot}}
\def\Hilbfunctor{\mathrm{Hilb}}
\def\Curvesstack{\mathcal{C}\!\mathit{urves}}
\def\Polarizedstack{\mathcal{P}\!\mathit{olarized}}
\def\Complexesstack{\mathcal{C}\!\mathit{omplexes}}
% \Pic is the operator that assigns to X its picard group, usage \Pic(X)
% \Picardstack_{X/B} denotes the Picard stack of X over B
% \Picardfunctor_{X/B} denotes the Picard functor of X over B
\def\Pic{\mathop{\mathrm{Pic}}\nolimits}
\def\Picardstack{\mathcal{P}\!\mathit{ic}}
\def\Picardfunctor{\mathrm{Pic}}
\def\Deformationcategory{\mathcal{D}\!\mathit{ef}}
\usepackage{cleveref}
\crefname{originalreference}{original reference}{original references}
\Crefname{originalreference}{Original reference}{Original references}
\makeatletter
\newlabel{lemma-duality-dim-1}{{0BS2}{source}{}{sourcepointer-1}{}}
\newlabel{lemma-duality-dim-1@cref}{{[originalreference][0][]0BS2}{[1][0][]source}{}{}{}}
\makeatother
\crefname{originalequation}{original equation}{original equations}
\Crefname{originalequation}{Original equation}{Original equations}
\makeatletter
\newlabel{lemma-equation-plane-curve}{{0BYB}{source}{}{sourcepointer-2}{}}
\newlabel{lemma-equation-plane-curve@cref}{{[originalequation][0][]0BYB}{[1][0][]source}{}{}{}}
\makeatother
\crefname{originalreference}{original reference}{original references}
\Crefname{originalreference}{Original reference}{Original references}
\makeatletter
\newlabel{lemma-genus-gorenstein}{{0C19}{source}{}{sourcepointer-3}{}}
\newlabel{lemma-genus-gorenstein@cref}{{[originalreference][0][]0C19}{[1][0][]source}{}{}{}}
\makeatother
\crefname{originalreference}{original reference}{original references}
\Crefname{originalreference}{Original reference}{Original references}
\makeatletter
\newlabel{lemma-genus-plane-curve}{{0BYD}{source}{}{sourcepointer-4}{}}
\newlabel{lemma-genus-plane-curve@cref}{{[originalreference][0][]0BYD}{[1][0][]source}{}{}{}}
\makeatother
\crefname{originalreference}{original reference}{original references}
\Crefname{originalreference}{Original reference}{Original references}
\makeatletter
\newlabel{lemma-genus-zero-positive-degree}{{0C6T}{source}{}{sourcepointer-5}{}}
\newlabel{lemma-genus-zero-positive-degree@cref}{{[originalreference][0][]0C6T}{[1][0][]source}{}{}{}}
\makeatother
\title{A Gorenstein proper curve of genus zero is a plane conic}
\author{The Stacks Project Authors}
\date{stacks-project:curves.tex:local-snapshot-20261012; GNU FDL 1.2 or later}
\begin{document}
\maketitle
\noindent\textit{Editorial note.} All delivered mathematical text is the author's own, copied from the pinned Stacks Project chapter file curves.tex; no statement or proof step was rewritten, reordered or removed.\par\medskip
\noindent\textit{Editorial note.} The counted proof is the author's complete argument: the degree computation, the closed immersion, the genus restriction to d in \{1, 2\}, and the linear-independence exclusion are all retained verbatim.\par\medskip
\noindent\textit{Editorial note.} The label lemma-genus-zero names the counted statement, and the proof environment immediately following it in the source is delivered verbatim as its proof.\par\medskip
\noindent\textit{Editorial note.} The counted statement is carried by the author's own lemma environment in the chapter source.\par\medskip
\noindent\textit{Allowed standard prerequisites (editorial).}\begin{itemize}
\item Degree of the dualizing sheaf in genus zero. Same chapter, Lemma lemma-genus-gorenstein. Gives deg(omega\_X) = -2, hence deg(omega\_X\textasciicircum{}tensor -1) = 2.
\item Closed immersion from a degree-2 invertible sheaf. Same chapter, Lemma lemma-genus-zero-positive-degree. Turns a basis of the three-dimensional space of global sections of L = omega\_X\textasciicircum{}tensor -1 into a closed immersion into the projective plane.
\item Plane-curve arithmetic genus. Same chapter, Lemmas lemma-equation-plane-curve and lemma-genus-plane-curve. Writes the equation F of the plane curve of degree d and forces d in \{1, 2\} from the vanishing genus.
\end{itemize}\medskip
\ifdefined\hypertarget\hypertarget{sourcepointer-1}{}\fi
\noindent\textit{Source reference.} Original reference 0BS2: sources/stacks-curves/pinned-original-source.tex, source lines 796--796. Prerequisite pointer invoked at proof line(s) 2116; it names a standard fact of the same chapter that the counted argument uses, and every step of the counted proof built on it is retained verbatim in the delivered body.\par\medskip
\ifdefined\hypertarget\hypertarget{sourcepointer-2}{}\fi
\noindent\textit{Source reference.} Original equation 0BYB: sources/stacks-curves/pinned-original-source.tex, source lines 1859--1859. Prerequisite pointer invoked at proof line(s) 2127; it names a standard fact of the same chapter that the counted argument uses, and every step of the counted proof built on it is retained verbatim in the delivered body.\par\medskip
\ifdefined\hypertarget\hypertarget{sourcepointer-3}{}\fi
\noindent\textit{Source reference.} Original reference 0C19: sources/stacks-curves/pinned-original-source.tex, source lines 1744--1744. Prerequisite pointer invoked at proof line(s) 2117; it names a standard fact of the same chapter that the counted argument uses, and every step of the counted proof built on it is retained verbatim in the delivered body.\par\medskip
\ifdefined\hypertarget\hypertarget{sourcepointer-4}{}\fi
\noindent\textit{Source reference.} Original reference 0BYD: sources/stacks-curves/pinned-original-source.tex, source lines 1904--1904. Prerequisite pointer invoked at proof line(s) 2129; it names a standard fact of the same chapter that the counted argument uses, and every step of the counted proof built on it is retained verbatim in the delivered body.\par\medskip
\ifdefined\hypertarget\hypertarget{sourcepointer-5}{}\fi
\noindent\textit{Source reference.} Original reference 0C6T: sources/stacks-curves/pinned-original-source.tex, source lines 2019--2019. Prerequisite pointer invoked at proof line(s) 2119; it names a standard fact of the same chapter that the counted argument uses, and every step of the counted proof built on it is retained verbatim in the delivered body.\par\medskip
\begin{lemma}
\label{lemma-genus-zero}
Let $X$ be a proper curve over a field $k$ with $H^0(X, \mathcal{O}_X) = k$.
If $X$ is Gorenstein and has genus $0$, then $X$
is isomorphic to a plane curve of degree $2$.
\end{lemma}

\begin{proof}
Consider the invertible sheaf $\mathcal{L} = \omega_X^{\otimes -1}$ where
$\omega_X$ is as in Lemma \ref{lemma-duality-dim-1}. Then
$\deg(\omega_X) = -2$ by Lemma \ref{lemma-genus-gorenstein}
and hence $\deg(\mathcal{L}) = 2$. By
Lemma \ref{lemma-genus-zero-positive-degree}
we conclude that choosing a basis $s_0, s_1, s_2$ of the $k$-vector
space of global sections of $\mathcal{L}$ we obtain a closed immersion
$$
\varphi_{(\mathcal{L}, (s_0, s_1, s_2))} :
X \longrightarrow \mathbf{P}^2_k
$$
Thus $X$ is a plane curve of some degree $d$. Let $F \in k[T_0, T_1, T_2]_d$
be its equation (Lemma \ref{lemma-equation-plane-curve}).
Because the genus of $X$ is $0$ we see that $d$ is $1$ or $2$
(Lemma \ref{lemma-genus-plane-curve}). Observe that
$F$ restricts to the zero section on $\varphi(X)$ and hence
$F(s_0, s_1, s_2)$ is the zero section of $\mathcal{L}^{\otimes 2}$.
Because $s_0, s_1, s_2$ are linearly independent we see that $F$
cannot be linear, i.e., $d = \deg(F) \geq 2$. Thus $d = 2$
and the proof is complete.
\end{proof}

\end{document}
